\section{A mixed surface with no index-zero ample class}
\label{sec:no-zero-class}

This section proves that the class-existence hypothesis in the ample-difference theorem fails on some mixed surfaces, and computes every ample polarization index on an explicit isogeny construction.

The construction starts from a geometric generic Jacobian with full symplectic monodromy. The passage from this monodromy to endomorphism ring $\mathbf Z$ is classical; compare \citet[Corollary~4.2, p.~297]{Chai}. Quotients of Picard-rank-one principally polarized varieties already supply classical obstructions to principal polarizations \citep[\S11.24, p.~169]{EGM}. We give the integral ordinary-index calculation, which is the feature needed here.

\begin{lemma}\label{lem:mixed-principal-rank-one}
There is an algebraically closed field $k$ of characteristic two and a principally polarized mixed abelian surface $(B,\Theta)/k$ such that
\[
             \End(B)=\mathbf Z,\qquad \NS(B)=\mathbf Z[\Theta].
\]
\end{lemma}
\begin{proof}
The smooth projective Artin--Schreier curve with function field
\[
                         y^2+y=x^3+\frac1x
\]
has genus two and $2$-rank one. Indeed its two poles have odd orders three and one. Their different contributions are four and two, so Riemann--Hurwitz gives $2g-2=-4+4+2=2$; the two-pole $2$-rank formula gives $f=2-1=1$. These formulas, with their hypotheses and proofs, are \citet[Lemmas~2.6--2.7 and Example~2.8, pp.~711--712]{PriesZhu}. In particular the $2$-rank-one stratum of the genus-two moduli space is nonempty.

Choose an irreducible component of this stratum over $\overline{\FF}_2$. Pass to a connected component with full level-five structure and then to a connected normal open scheme carrying the universal curve and its principally polarized Jacobian. The geometric monodromy theorem of \citet[Theorem~4.5, pp.~14--15 of the author version; scheme rigidifications in \S5.1, pp.~17--18]{AchterPries} applies in characteristic two, genus two and $2$-rank one. Before the finite level cover its $3$-adic image is $\operatorname{Sp}_4(\mathbf Z_3)$. A finite cover changes this to a finite-index subgroup; restriction to a nonempty open does not decrease the image, since a connected finite \'etale cover of a normal integral scheme remains connected over that open. Its Zariski closure is therefore $\operatorname{Sp}_{4,\mathbf Q_3}$. Let $B$ be the Jacobian at the geometric generic point, over an algebraic closure $k$ of the function field. Its canonical principal polarization is represented by an ample line bundle $\Theta$, and $B$ is mixed.

We spell out why the geometric endomorphism ring has rank one. It is a finite free abelian group, and all its generators are defined over a common finite extension of the generic field \citep[Theorem~12.10 and Corollaries~12.11--12.13, pp.~182--184]{EGM}. For the finite-generation step, the classical positive integral form on the endomorphism group gives an integer Gram matrix on a rational basis. Its nonzero determinant bounds all denominators of every endomorphism, so the group lies in a fixed finite-rank lattice. After spreading out and shrinking, the generators act on the abelian scheme over the normalization in that extension. A purely inseparable extension leaves the \'etale fundamental group unchanged, and the separable part becomes finite \'etale after shrinking. Thus every generator commutes with a finite-index monodromy subgroup, still Zariski dense in $\operatorname{Sp}_4$.

The centralizer of $\operatorname{Sp}_4$ on its defining representation is $\mathbf Q_3$ times the identity. Commutation with the diagonal torus makes an endomorphism diagonal in the four standard weight spaces; the elementary symplectic root matrices force its four diagonal entries to agree. The Tate-module map is injective after tensoring $\End(B)$ with $\mathbf Z_3$: an endomorphism killing $B[3^m]$ factors through the quotient $[3^m]:B\to B$, so its kernel modulo $3^m$ is exactly $3^m\End(B)$. Taking inverse limits, using finite generation, gives the injection; this is also \citet[Lemmas~12.8--12.9, pp.~182--183]{EGM}. Hence
\[
      \End(B)\otimes\mathbf Q_3\ \subset\ \mathbf Q_3\operatorname{id}.
\]
It follows that $\End^0(B)=\mathbf Q$. A finite $\mathbf Z$-algebra contained in $\mathbf Q$ consists of rational algebraic integers, so $\End(B)=\mathbf Z$.

Finally $L\mapsto\lambda_\Theta^{-1}\lambda_L$ injects $\NS(B)$ into $\End(B)$: the kernel on $\Pic(B)$ is exactly $\Pic^0(B)$ \citep[\S\S7.22--7.26, pp.~106--108]{EGM}. It sends $[\Theta]$ to the identity. Therefore $\NS(B)=\mathbf Z[\Theta]$.
\end{proof}

\begin{proof}[Proof of Theorem~\ref{thm:no-ordinary-index-zero-ample-class}]
Choose $(B,\Theta)$ as in the lemma. The local-local height-two summand of $B[2]$ contains a subgroup $H\simeq\alpha_2$. More explicitly, the mixed surface has
\[
 B[2]\simeq(\mathbf Z/2\mathbf Z\oplus\mu_2)\oplus I_{1,1},
\]
where $I_{1,1}$ is the $2$-torsion of a supersingular elliptic curve and contains a unique $\alpha_2$; see \citet[Example~2.2, Lemma~3.1 and Propositions~3.2--3.3, author version pp.~1--4; genus-two table p.~6]{PriesTorsion}. Naturality of relative Frobenius gives $H\subset\ker(F_B:B\to B^{(2)})$.

Form the quotient isogeny $q:B\to A=B/H$. Its degree is two and it is radicial. The relative Frobenius, of degree four, factors as
\[
                    F_B=hq,\qquad h:A\longrightarrow B^{(2)},
\]
where $h$ is an isogeny of degree two. Define the actual line bundle
\begin{equation}\label{eq:no-zero-pullback}
 L=h^*(\Theta^{(2)}),\qquad q^*L\simeq F_B^*(\Theta^{(2)})
                              \simeq\Theta^{\otimes2}.
\end{equation}
It is ample because $h$ is finite. The isogeny preserves the $2$-rank, so $A$ is mixed. Riemann--Roch and pullback by the degree-two isogeny give
\[
                       2\chi(L)=\chi(2\Theta)=4,
                       \qquad \chi(L)=2.
\]

The kernels of both $q$ and $q^\vee$ are $\alpha_2$, since this group is Cartier self-dual. Their maps on integral geometric $2$-adic Tate modules are isomorphisms. Indeed a radicial isogeny is bijective on algebraically closed points; if its image of $x$ is killed by $2^m$, then $2^mx$ is a point of its infinitesimal kernel and must be zero. Thus the bijection restricts to every geometric $2^m$-torsion group. The polarization identity from \eqref{eq:no-zero-pullback} is
\[
                         q^\vee\lambda_Lq=2\lambda_\Theta.
\]
The outer Tate-module maps are integral isomorphisms and $\lambda_\Theta$ is principal. On the rank-one Tate modules the middle map therefore has cokernel length one. Consequently
\[
                         a(L)=1,\qquad b(L)=v_2(2)-1=0.
\]

Pullback by $q$ injects N\'eron--Severi groups: its kernel is killed by the degree through the norm identity, whereas $\NS(A)$ is torsion-free by its injection into $\Hom(A,A^\vee)$. For any line bundle $N$ on $A$, write $[q^*N]=m[\Theta]$. Euler characteristic depends only on this class, and
\[
                            2\chi(N)=m^2.
\]
Thus $m$ is even. Equation~\eqref{eq:no-zero-pullback} and injectivity imply $[N]=(m/2)[L]$. This proves $\NS(A)=\mathbf Z[L]$. If $N$ is ample then $t=m/2>0$, because $q^*N$ is ample. Algebraically trivial twists do not change the polarization homomorphism, so
\[
 \lambda_N=t\lambda_L,\qquad \chi(N)=2t^2,
 \qquad a(N)=1+v_2(t),\qquad b(N)=v_2(t).
\]
In particular no ample $N$ has ordinary index zero. Theorem~\ref{thm:finite-field} also shows that this $A$ cannot have a finite-field model.
\end{proof}

The same descent follows from the following theta-group calculation, an elementary special case of the classical finite-group extension statement \citep[Fact~8.7, p.~115]{EGM}.

\begin{lemma}[$\alpha_2$ theta descent]\label{lem:alpha-two-theta-descent}
Over an algebraically closed field of characteristic two, every central extension of fppf group schemes
\[
                  1\longrightarrow\mathbf G_m\longrightarrow E
                   \longrightarrow\alpha_2\longrightarrow1
\]
splits. If $H\simeq\alpha_2$ is a subgroup of the polarization kernel $K(P)$ of a line bundle $P$ on an abelian variety, then $P$ descends to the quotient by $H$.
\end{lemma}
\begin{proof}
The underlying $\mathbf G_m$-torsor is trivial because $\alpha_2=\operatorname{Spec}k[t]/(t^2)$ has local coordinate ring and hence trivial Picard group. Choose a section taking zero to the identity. Centrality makes multiplication on the resulting torsor take the form
\[
        (u,x)(v,y)=(uvc(x,y),x+y),\qquad
 c\in\bigl(k[x,y]/(x^2,y^2)\bigr)^\times.
\]
The identity conditions give $c(x,0)=c(0,y)=1$, so $c(x,y)=1+\gamma xy$ for some $\gamma\in k$. In particular it is symmetric. Choose $\beta^2=\gamma$ and set $d(t)=1+\beta t$. Direct multiplication in this coordinate ring gives
\[
                \frac{d(x)d(y)}{d(x+y)}=1+\beta^2xy=c(x,y),
\]
since $(x+y)^2=x^2=y^2=0$. Multiplying the chosen section by $d(t)^{-1}$ makes its cocycle one. It is then a group homomorphism and splits the extension.

Restrict the theta-group extension of $P$ to $H$. Its splitting is precisely an $H$-linearization of $P$. Faithfully flat descent along the quotient torsor gives the required line bundle; see \citet[Theorem~8.10 and its proof, p.~116]{EGM}. In the construction above $H\subset B[2]=K(2\Theta)$, so this also descends $\Theta^{\otimes2}$ and yields the same polarization-index calculation.
\end{proof}
